\documentclass[10pt,a4paper]{amsart}
\usepackage[margin=19mm]{geometry}
\usepackage{amsmath,amssymb,mathtools}
\usepackage[scr=rsfs,cal=euler]{mathalfa}
\usepackage{xfrac}
\usepackage[unicode,hidelinks,bookmarks=false]{hyperref}
\newcommand{\ie}{i.e.\ }
\newcommand{\cf}{cf.\ }
\newcommand{\resp}{respectively\ }
\newcommand{\Subsection}{\subsection}
\providecommand{\nobreakdash}{\nobreak\hbox{-}}
\setlength{\emergencystretch}{2em}
\multlinegap0pt
\usepackage{bm}
\usepackage{bbm}
\usepackage{dsfont}
\usepackage{xspace}
\usepackage{cancel}
\usepackage{mathtools}
\usepackage{amsthm}
\makeatletter
\@ifundefined{theorem}{\newtheorem{theorem}{Theorem}}{}
\@ifundefined{definition}{\newtheorem{definition}[theorem]{Definition}}{}
\makeatother
\makeatletter
\newcommand{\B}{\mathbf{B}}
\newcommand{\Ss}{\mathbf{S}}
\expandafter\ifx\csname acknowledgement\endcsname\relax
\newenvironment{acknowledgement}{\noindent\textbf{Acknowledgments}}{}
\fi
\expandafter\ifx\csname proof\endcsname\relax
\newenvironment{proof}[1][Proof]{\noindent\textbf{#1.} }{\hfill\rule{0.5em}{0.5em}}
\fi
\makeatother
\title[Moduli space theory for complete, constant Q-curvature metri]{Moduli space theory for complete, constant Q-curvature metrics on finitely punctured spheres}
\author{Rayssa Caju, Jesse Ratzkin, Almir Silva Santos}
\date{Source: arXiv:2403.06511v1 (CC BY 4.0)}
\begin{document}
\maketitle

\noindent\emph{Source and licence.} Moduli space theory for complete, constant Q-curvature metrics on finitely punctured spheres. Version arXiv:2403.06511v1 (CC BY 4.0), distributed under the Creative Commons Attribution 4.0 International licence, as recorded in the arXiv licence metadata for this exact version and shown on the abstract page https://arxiv.org/abs/2403.06511v1. Full rights notice: licenses/arxiv-2403.06511-CC-BY-4.0.txt.\par\smallskip

\noindent\textit{Editorial scope.} Counted unit: the Theorem whose source label is \texttt{thm:symp\_form} in the article (source0.tex lines 1165--1168, source label \texttt{thm:symp\_form}), delivered together with the article's own complete original proof of it (source0.tex lines 1170--1303, one \texttt{proof} environment of 134 lines). No mathematics is written, repaired or re-derived: statement and proof are the article's own text line for line. Required context retained uncounted: lin\_op\_eucl (lines 434--436); emden\_fowler\_coords (lines 463--467); del\_ode (lines 499--502); eq002 (lines 505--509); geom\_jac\_fields (lines 663--665); deficiency\_defn (lines 773--784); defn\_symp\_form (lines 1158--1160). Every definition, display and label of those lines is retained unchanged. Omitted from this package and not redistributed: 1--433, 437--462, 468--498, 503--504, 510--662, 666--772, 785--1157, 1161--1164, 1169--1169, 1304--1325. Nothing inside a retained excerpt is omitted or altered: every retained body is delivered exactly as the article's lines are written, so no omission receipt of any kind accompanies this package. Citations: the retained excerpts carry no citation command at all, so no bibliography entry of the article is transcribed and no reference is redistributed. Reference adaptation: none was needed and none was made; every reference inside the delivered unit resolves inside it, so no unresolved reference or citation is produced. Printed slips kept literal: no printed word, symbol, hypothesis or number of the counted statement or of its proof was altered. Layout and class: the article's own class is \texttt{amsart} with packages including amsmath, amssymb, latexsym, soul, cite, mathrsfs, color, enumitem, graphicx, hyperref, babel, geometry, backref; the wrapper is the lane's pinned \texttt{amsart} editorial wrapper at 10pt on A4, which restates the theorem-like environments the retained text needs and adds the article's own macro definitions that the retained text needs (\texttt{\textbackslash B}, \texttt{\textbackslash Ss}), with extra wrapper packages (bm, bbm, dsfont, xspace, cancel, mathtools). The delivered numbering is the unit's own. Assets: no retained span contains a figure environment, a graphics-inclusion command, a PDF-page inclusion, a table or a TikZ picture; no image file is copied into this package, the package declares no asset dependency and no image file is redistributed with it. Licence: version arXiv:2403.06511v1 of this article is distributed under the Creative Commons Attribution 4.0 International licence, as recorded in the arXiv licence metadata for this exact version and shown on the abstract page https://arxiv.org/abs/2403.06511v1. A full rights notice is delivered at licenses/arxiv-2403.06511-CC-BY-4.0.txt. Characters: every retained excerpt is pure ASCII, so the delivered engine, XeLaTeX with its default Latin Modern text font, reports no missing character.
\begin{equation} \label{lin_op_eucl} 
L_g(v)  =  (-\Delta_0)^2 v - \frac{n(n+4)(n^2-4)}{16} u^{\frac{8}{n-4}} v. 
\end{equation} 

\begin{equation} \label{emden_fowler_coords} 
\mathfrak{F} : \mathcal{C}^\infty (\B_r(0) \backslash \{ 0 \}) 
\rightarrow \mathcal{C}^\infty ((-\log r, \infty) \times \Ss^{n-1}), 
\qquad \mathfrak{F}(u) (t,\theta) = e^{\frac{4-n}{2}t} u(e^{-t}\theta). 
\end{equation} 

\begin{equation} \label{del_ode} 
\ddddot v -  \frac{n(n-4)+8}{2} \ddot v + \frac{n^2(n-4)^2}{16} 
v - \frac{n(n-4)(n^2-4)}{16} v^{\frac{n+4}{n-4}}=0. 
\end{equation} 

\begin{equation}\label{eq002}
    \mathcal{H}_\varepsilon = - \dot v_\varepsilon\dddot v_\varepsilon + \frac{1}{2} \ddot v_\varepsilon^2
+  \frac{n(n-4)+8}{4} \dot v_\varepsilon^2- \frac{n^2(n-4)^2}{32} 
v_\varepsilon^2 + \frac{(n-4)^2(n^2-4)}{32} v_\varepsilon^{\frac{2n}{n-4}}.
\end{equation}

\begin{equation} \label{geom_jac_fields}
w_0^+(\varepsilon) = \dot v_{\varepsilon} , \qquad w_0^-(\varepsilon) 
= \frac{d}{d\varepsilon}v_\varepsilon. \end{equation} 

\begin{definition} \label{deficiency_defn}
Let $g\in \mathcal{M}_\Lambda$ and choose $r_0>0$ sufficiently 
small such that $\mathbf{B}_{2r_0} (p_i) \cap \mathbf{B}_{2r_0}(p_j) 
= \varnothing$ for each distinct pair of punctures. We define the 
deficiency space $\mathcal{W}_g$ by 
$$\mathcal{W}_g = \operatorname{Span} \{ \chi \mathfrak{F}^{-1} (w_0^+(\varepsilon_i)), 
\chi \mathfrak{F}^{-1} (w_0^-(\varepsilon_i)): i=1, \dots, k\},$$
where $\chi$ is a fixed cut-off function such that 
$$\chi(x) = \left \{ \begin{array}{rl} 1 & |x| < r_0 \\ 0 & |x|> 3r_0/2 
\end{array} \right . , \qquad \|\nabla^k \chi \|_{\mathcal{C}^0} 
\leq c r^{-k}. $$
\end{definition}

\begin{equation} \label{defn_symp_form}
\omega(v,w) = \lim_{r \searrow 0} \int_{\Omega_r} (v L_g (w) - 
w L_g (v) )d\mu_0.\end{equation} 

\begin{theorem} \label{thm:symp_form} 
The form $\omega$ defined in \eqref{defn_symp_form} is a 
symplectic form on the $2k$-dimensional vector space $\mathcal{W}_g$. 
\end{theorem}

\begin{proof} Our first order of business is to show that 
$\omega$ is well-defined, {\it i.e.} that the limit in \eqref{defn_symp_form}
exists. By \eqref{lin_op_eucl} observe that 
$$vL_g(w) - wL_g(v) = v \Delta_0^2 w - w \Delta_0^2 v.$$
Next, we recall that the outer unit normal of $\Omega_r$ is 
$-\partial_r$ on each boundary sphere $\partial \B_r(p_i)$ 
and integrate by parts to see 
$$\int_{\Omega_r} v\Delta_0^2 w d\mu_0 = \int_{\Omega_r} 
(\Delta_0 v) (\Delta_0 w) d\mu_0 - \sum_{i=1}^k \int_{\partial 
\B_r(p_i)} (v \partial_r \Delta_0 w - \partial_r v \Delta_0 w
)d\sigma_0,$$
and so \eqref{defn_symp_form} becomes 
 \begin{equation} \label{symp_form2} 
 \omega(v,w) = \lim_{r \searrow 0} \sum_{i=1}^k 
 \int_{\partial \B_r(p_i)} (w \partial_r \Delta_0 v - v \partial_r \Delta_0 w
 + \partial_r v \Delta_0 w - \partial_r w \Delta_0 v )d\sigma_0.\end{equation} 

 Next, we change variables using \eqref{emden_fowler_coords}, letting 
 $\widetilde v = \mathfrak{F}(v)$ and $\widetilde w= \mathfrak{F}(w)$. Under this 
 change of variables 
 \begin{equation} \label{transformed_derivatives1} 
 \partial_r v = -e^{ \frac{n-2}{2}t} 
 \left ( \partial_t \widetilde v + \frac{n-4}{2} \widetilde v \right ), 
 \qquad \Delta_0 v = e^{\frac{n}{2}t} \left ( \partial_t^2 \widetilde v - 2 \partial_t 
 \widetilde v -\frac{n(n-4)}{4} \widetilde v + \Delta_\theta \widetilde v \right )
 \end{equation} 
 and 
 \begin{equation} \label{transformed_derivatives2} 
 \partial_r \Delta_0 v = -e^{\frac{n+2}{2}t} \left ( \partial_t^3 
 \widetilde v + \frac{n-4}{2} \partial_t^2 \widetilde v -\frac{n^2}{4} \partial_t \widetilde v
 -\frac{n^2(n-4)}{8} \widetilde v + \Delta_\theta \partial_t \widetilde v + \frac{n}{2} \Delta_\theta 
 \widetilde v \right ).\end{equation} 


 Plugging \eqref{transformed_derivatives1} and \eqref{transformed_derivatives2}
 into \eqref{symp_form2} we obtain 
 \begin{align*}
   \int_{\partial \B_r(p_i)} & \left(w \partial_r \Delta_0 v - v \partial_r \Delta_0 w + 
 \partial_r v \Delta_0 w - \partial_r w \Delta_0 v\right) d\sigma_0\\  
 = & \int_{\Ss^{n-1}} (w \partial_r \Delta_0 v - v \Delta_0 w + \partial_r v \Delta_0 w - \partial_r w  \Delta_0 v )(r\theta)r^{n-1} d\theta\\
 = &\int_{\Ss^{n-1}} \left ( -\widetilde{w}\partial_t^3 \widetilde{v} - \frac{n-4}{2}  \widetilde{w}
 \partial_t^2 \widetilde{v} + \frac{n^2}{4} \widetilde{w}\partial_t \widetilde{v} + \frac{n^2(n-4)}{8} \widetilde{w}\widetilde{v} 
 - \widetilde{w}\Delta_\theta \partial_t \widetilde{v} - \frac{n}{2} \widetilde{w}\Delta_\theta \widetilde{v}  \right.\\
 & +\widetilde{v}\partial_t^3 \widetilde{w} + \frac{n-4}{2}  \widetilde{v}
 \partial_t^2 \widetilde{w} - \frac{n^2}{4} \widetilde{v}\partial_t \widetilde{w} - \frac{n^2(n-4)}{8} \widetilde{v}\widetilde{w} 
 + \widetilde{v}\Delta_\theta \partial_t \widetilde{w} + \frac{n}{2} \widetilde{v}\Delta_\theta \widetilde{w} \\ 
 &-\partial_t \widetilde{v} \partial_t^2 \widetilde{w} + 2 \partial_t \widetilde{v} 
 \partial_t \widetilde{w} + \frac{n(n-4)}{4} \widetilde{w} \partial_t \widetilde{v} - \partial_t \widetilde{v} 
 \Delta_\theta \widetilde{w}  - \frac{n-4}{2} \widetilde{v} \partial_t^2\widetilde{w} + (n-4)\widetilde{v} 
 \partial_t\widetilde{w}\\
 &+ \frac{n(n-4)^2}{8} \widetilde{v}\widetilde{w} - \frac{n-4}{2} \widetilde{v} \Delta_\theta
 \widetilde{w}\\ 
 &+\partial_t \widetilde{w} \partial_t^2 \widetilde{v} - 2 \partial_t \widetilde{w} 
 \partial_t \widetilde{v} - \frac{n(n-4)}{4} \widetilde{v} \partial_t \widetilde{w} + \partial_t \widetilde{w} 
 \Delta_\theta \widetilde{v}  + \frac{n-4}{2} \widetilde{w} \partial_t^2\widetilde{v} - (n-4)\widetilde{w} 
 \partial_t\widetilde{v}\\
 &\left.- \frac{n(n-4)^2}{8} \widetilde{w}\widetilde{v} + \frac{n-4}{2} \widetilde{w} \Delta_\theta
 \widetilde{v} \right)d\theta\\
 = & \int_{\Ss^{n-1}} \left ( \widetilde{v} \partial_t^3 \widetilde{w} - \widetilde{w} \partial_t^3 \widetilde{v} 
 + \partial_t\widetilde{w} \partial_t^2 \widetilde{v} - \partial_t{v} \partial_t^2\widetilde{w} 
 +\frac{n(n-4)+8}{2} (\widetilde{w}\partial_t \widetilde{v} - \widetilde{v} \partial_t\widetilde{w} 
 )\right ) d\theta
 \end{align*}


 Observe that each element of the deficiency space is asymptotically radial about 
 each puncture point, so that the expansions 
 $$\mathfrak{F} (v(\cdot - p_i)) (t,\theta) = \overline{v}(t) + \mathcal{O} (e^{-\delta t}), 
 \qquad \mathfrak{F}(w(\cdot - p_i)) (t,\theta) = \overline{w}(t) + \mathcal{O}(e^{-\delta t})$$
 for each puncture $p_i$. Thus each term in the expansion above involving derivatives with 
 respect to $\theta$ will vanish in the limit. 

 Next we use the fact that both $v$ and $w$ lie in the definciency space 
 $\mathcal{W}_g$ (see Definition \ref{deficiency_defn}). This means that near each puncture $p_i$ the 
 functions $v$ and $w$ have asymptotic expansions of the form 
 \begin{equation} \label{expansion_def_space}
 \widetilde v = \alpha_i^+ w_0^+(\varepsilon_i) + \alpha_i^- w_0^-(\varepsilon_i), \qquad 
 \widetilde w = \beta_i^+ w_0^+(\varepsilon_i) + \beta_i^- w_0^-(\varepsilon_i)  
 \end{equation} 
 for $t$ sufficiently large. 


 At each end, using bilinearity and skew-symmetry, we see that $\omega(v,w)$ is the limit as $t\rightarrow\infty$ of
\begin{eqnarray}\label{eq001}
    (\alpha_i^+\beta_i^--\alpha_i^-\beta_i^+)\int_{\Ss^{n-1}}\left(w_0^+ \dddot {w}_0^- 
     -  w_0^- \dddot {w}_0^+ + \dot {w}_0^- \ddot {w}_0^+ - \dot {w}_0^+ \ddot {w}_0^-
     \right .&& \\ \nonumber 
     \qquad \qquad \qquad \qquad \qquad  \left .  + \frac{n(n-4)+8}{2} (w_0^- \dot {w}_0^+ - w_0^+ \dot {w}_0^-)\right)d\theta & & 
\end{eqnarray}

If $A_\varepsilon$ is the integrand in \eqref{eq001}, then
\begin{align*}
\frac{d}{dt}A_\varepsilon
& = w_0^+ \left(\ddddot w_0^--\frac{n(n-4)+8}{2}\ddot w_0^-\right) - w_0^- \left(\ddddot w_0^+-\frac{n(n-4)+8}{2}\ddot w_0^+\right)\\
& =\left ( 
\frac{n^2(n-4)^2}{16} - \frac{n(n+4)(n^2-4)}{16} v_\varepsilon^{\frac{8}{n-4}} \right )\left(w_0^+w_0^--w_0^-w_0^+\right)=0.
\end{align*}

Thus $A_\varepsilon$ does not depend on $t$. Here we have used the ODE for $w_0^\pm$, namely 
$$\ddddot w_0^\pm - \frac{n(n-4)+8}{2} \ddot w_0^\pm + \left ( 
\frac{n^2(n-4)^2}{16} - \frac{n(n+4)(n^2-4)}{16} v_\varepsilon^{\frac{8}{n-4}} \right ) w_0^\pm 
= 0.$$

Let us find the value of the integrand in \eqref{eq001} at $t=0$ using the definitions of $w_0^\pm$ in \eqref{geom_jac_fields}. First observe that $w_0^+(\varepsilon) = \dot v_\varepsilon$, 
so $w_0^+ (0)  =  \dot v_\varepsilon (0) = 0$, $\dot w_0^+ (0)  =  \ddot v_\varepsilon (0) > 0$, $\ddot w_0^+(0)  =  \dddot v_\varepsilon (0) = 0$ and by \eqref{del_ode} we get
$$\dddot w_0^+ (0) = \ddddot v_\varepsilon (0)  =  \frac{n(n-4)+8}{2} \ddot v_\varepsilon(0) - \frac{n(n-4)}{16} 
\left ( n(n-4) \varepsilon - (n^2-4)\varepsilon^{\frac{n+4}{n-4}} \right ) .
$$
Furthermore, since $\displaystyle w_0^- = \frac{d}{d\varepsilon} v_\varepsilon$
and $v_\varepsilon$ assumes its minimal value at $t=0$ we see 
$w_0^-(0) =1$. Therefore, at $t=0$ it holds
\begin{align*}
    A_\varepsilon(0) & =-\dddot w_0^+(0)-\ddot v_\varepsilon(0)\ddot w_0^-+\frac{n(n-4)+8}{2}\ddot v_0^+(0)\\
    & =-\ddot v_\varepsilon(0)\ddot w_0^-+\frac{n(n-4)}{16} 
\left ( n(n-4) \varepsilon - (n^2-4)\varepsilon^{\frac{n+4}{n-4}} \right )
\end{align*}
By \eqref{eq002} we have
\begin{equation} \label{eq003}
\mathcal{H}_\varepsilon = \frac{1}{2} \ddot v_\varepsilon (0)^2 - \frac{n^2(n-4)^2}{32} 
\varepsilon^2 + \frac{(n-4)^2(n^2-4)}{32} \varepsilon^{\frac{2n}{n-4}} . 
\end{equation} 
As we remarked earlier, the energy is a decreasing function of $\varepsilon$, minimized by the 
cylinder, which has the largest possible necksize, so differentiating \eqref{eq003} we see 
\begin{align}
    0  >  \frac{d}{d\varepsilon} \mathcal{H}_\varepsilon & = \ddot v_\varepsilon (0) 
\frac{d}{d\varepsilon} \ddot v_\varepsilon (0) - \frac{n^2(n-4)^2}{16} \varepsilon 
+ \frac{n(n-4)(n^2-4)}{16} \varepsilon^{\frac{n+4}{n-4}} \\ \nonumber 
 & =  \ddot v_\varepsilon (0) \ddot w_0^- (0) - \frac{n(n-4)}{16}(n(n-4) \varepsilon 
- (n^2-4)\varepsilon^{\frac{n+4}{n-4}} ) .
\end{align}
This implies that
$$A_\varepsilon(0)=-\frac{d}{d\varepsilon}\mathcal H_\varepsilon>0,$$
and so $\omega$ is nondegenerate.
 \end{proof}

\end{document}
